\summarychapter{3}{Only Worthless People}{Catholic Christianity}

\begin{inbrief}
Between the death of the apostles and the first Christian emperor the scattered congregations grew
into one body with a universal vision. The chapter takes a ``flying tour'' of where the faith spread,
looks at the kind of people who joined it, and asks what made it grow.
\question{How did the scattered congregations of apostolic times become catholic Christianity---where
did the faith spread, and why so rapidly?}
\end{inbrief}

Eusebius, the early church historian, tells how King Abgar the Black of Edessa, capital of the small
kingdom of Osrhoene beyond the Roman border, wrote to Jesus asking to be healed, and how Jesus promised
to send a disciple after his ascension. The story is legend, but it is a reminder that the gospel did
not stop at the frontier: Osrhoene became the first Christian kingdom and a link to lands farther east.

First-century Christianity was a spiritual explosion set off by ``the Event,'' the presence of Jesus
Christ. The second and third centuries channeled that power, let the church come to terms with time
and gave it a lasting character: \emph{catholic} Christianity, more than an organization, the
conviction that all Christians should be one body. It carried on Jesus' commission and Paul's work.
The word is not in the New Testament. Ignatius, bishop of Antioch early in the second century, first
used it: ``Wherever Jesus Christ is, there is the catholic church.'' By the end of that century it
meant the church that is both universal, as against local congregations, and orthodox, as against
heretical groups. By the reign of Constantine (312--337) there were churches in every large town of
the empire and as far apart as Britain, Carthage and Persia.

\begin{sidebox}{Did Christianity Grow Miraculously?}
The Bible reports miracles that accompanied the church's growth, but the growth rate itself need not
have been miraculous. Theorists such as Rodney Stark hold that a steady rate, like that of Mormonism in
America, accounts for the numbers. Christianity's greatest expansion has come in the last fifty years.
\end{sidebox}

\ssection{The Spread of the Faith}
The tour begins with the Jews, who lived in every part of the empire and may have been 7 percent of its
population. Their faith attracted many Gentiles. A few accepted circumcision; most remained
``God-fearers'' on the edge of the synagogue. These people knew the Old Testament, its ideas and its
morals, and when Christian preachers offered them all that Judaism offered and more, without
circumcision, they responded readily. No later missionary movement found a field so well prepared.
Like the synagogues, the churches were local assemblies that saw themselves as one faithful Israel
across the world, which helps explain the catholic outlook.

The Roman world was a world of cities, and Paul had set the pattern: settle in a great city and reach
the surrounding towns through helpers. After AD 70 the center of the movement shifted north and then
west. Antioch in Syria became the church's second home; by the end of the fourth century half of its
half-million people were Christians. Edessa, tied to Antioch, later claimed Addai, one of the seventy
disciples, as its founder, and about 200 Serapion of Antioch consecrated Palut as its bishop. From
there some Christian probably went on to India. The ``Thomas Christians'' say it was the apostle
Thomas; that is possible but cannot be proved. What is fairly certain is that the church in India is
very old.

\ssection{The Move West}
The main stream flowed west, along Paul's course toward Italy and Spain. Ephesus and the Greek-speaking
cities of Asia Minor were fertile ground. Even rural Bithynia saw such growth that in 112 its governor
Pliny wrote to the emperor Trajan about Christians of every age, class and sex in town and country, and
feared the temples would be deserted---perhaps the first mass movement in Christian history. This was
unusual: backward regions with their own languages generally resisted, and in the sixth century
Justinian was still fighting paganism in inner Asia Minor.

In Rome the church, planted by unknown believers, grew fast. Adolf Harnack reckoned at least thirty
thousand Christians there by AD 250, most of them poor, since for over a century the Roman church spoke
Greek, the language of slaves, not the Latin of the upper class. With its claim to Peter and Paul and
its place in the capital, it won respect and a natural leadership. Farther west and north progress was
slower: a church at Lyons in Gaul by the mid-second century, known through its bishop Irenaeus;
bishops in Spain by the end of the third. How the faith reached Britain is unknown; three British
bishops attended the Council of Arles in 314.

\ssection{North Africa}
Across the Mediterranean the center was Carthage. Every town and almost every village had its bishop.
Here arose the first Latin-speaking churches, drawn largely from the Romanized upper class, while the
Punic and Berber speakers around them made for tensions of race and language that would cause trouble
in the third-century persecutions. Cyrene, to the east, appears four times in the New Testament---Simon
who carried the cross, the pilgrims at Pentecost, Stephen's opponents and the first evangelists
of Gentiles---and surely had churches early. Alexandria, founded by Alexander the Great in 332 BC,
was the empire's second city, with a large Jewish community whose philosopher Philo tried to explain
Judaism in the terms of Greek philosophy. Christians there faced the same task in their famous
catechetical school. They claimed John Mark as founder, and in the third and fourth centuries few
churches had more influence. In sum, by the end of the third century no part of the empire was without
a witness; it was strongest in Syria, Asia Minor, North Africa and Egypt, and in cities such as Rome
and Lyons. Little is known of the countryside.

\ssection{The Social Impact of the Gospel}
The catholic vision also crossed social lines. For three centuries most believers were humble people:
slaves, women, traders, soldiers. The critic Celsus sneered that Christians could convince ``only
worthless and contemptible people.'' Shelley counts it to the church's credit that it did not neglect
the poor; but by the end of the second century some of the keenest minds were also becoming Christians.
The \emph{apologists} (the word means defenders, as in a court) answered pagan slanders and exposed the
weakness of paganism, addressing emperors but aiming at the educated public: Aristides, Justin Martyr,
Tatian, Athenagoras, Theophilus of Antioch, the author of the \emph{Letter to Diognetus} and Melito of
Sardis. Irenaeus of Lyons wrote five books against the gnostic heresies and a \emph{Proof of the
Apostolic Preaching}; his biblical theology steadied the church and prepared later Christian views of
history such as Augustine's. Tertullian of Carthage, born about 150, the ``father of Latin theology,''
attacked the absurdity of persecution in his \emph{Apology}, encouraged martyrs, fought heretics and
was the first to use the Latin word \emph{trinitas}. In Alexandria the converted Stoic Pantaenus, who
probably traveled to India, was teaching by about 185, and his pupil Clement carried the school
further. By the third century the church, with its travel, synods, letters and loyalty, looked like an
empire within the empire.

\ssection{Reasons for the Gospel's Spread}
Believers stress the power of God; but God works through people, and Shelley names four human factors.
First, burning conviction: the Event had happened, and Christians could not keep the news of
redemption to themselves, even in the face of martyrdom. Second, the gospel met a felt need. Stoicism
sought tranquility by detachment and courage; Christians could admire its courage in suffering, its
independence of worldly things and its trust in providence, and many saw that what the Stoics aimed
at, the Holy Spirit produced. Third, practical love: care for the poor, widows, orphans, prisoners and
those condemned to the mines, and help in famine, earthquake and war. ``See how these Christians love
one another,'' pagans said. Churches buried poor members---Lactantius insisted that God's image must
not be left to beasts and birds---and from the later second century acquired burial grounds, such as
the one at Catacumbas on the Appian Way, which is how Christians came to be linked with the catacombs.
Even the apostate emperor Julian, trying to revive the old religion, admitted that the ``godless
Galileans'' cared for pagan poor as well as their own. Fourth, persecution advertised the faith:
``martyr'' means witness, and the calm courage, courtesy and joy of the martyrs, especially the young
women, moved the crowds and sometimes converted onlookers on the spot.

So the churches multiplied until Rome could neither ignore nor suppress them. The age shows, Shelley
concludes, that the church is truly catholic only when the gospel drives it to bring all people to
living faith in Christ.

\begin{keyfacts}
\keypeople{Eusebius; Abgar of Edessa; Ignatius of Antioch; Pliny and Trajan; Irenaeus of Lyons; Celsus;
Justin Martyr and the other apologists; Tertullian; Pantaenus; Clement of Alexandria; Philo;
Lactantius; Julian the Apostate; modern scholars Adolf Harnack, Rodney Stark, Ward Gasque.}
\keydates{%
  \kd{332 BC}{Alexander the Great founds Alexandria}
  \kd{c. 112}{Pliny writes to Trajan about the Christians of Bithynia}
  \kd{c. 150}{Birth of Tertullian at Carthage}
  \kd{c. 185}{Pantaenus teaching at Alexandria}
  \kd{c. 200}{Serapion of Antioch consecrates Palut bishop of Edessa}
  \kd{250}{At least 30,000 Christians in Rome (Harnack's estimate)}
  \kd{314}{Three British bishops attend the Council of Arles}}
\keyterms{Catholic (universal and orthodox); God-fearers; apologists; \emph{trinitas}; catechetical
school; catacombs; martyr (``witness''); Stoicism.}
\end{keyfacts}
